\section{Details of Non-Implications}
\label{sec:cex-extra}

\subsection{Trivial Examples}
\label{sec:cex-extra:trivial}

\begin{example}[single item]
\label{cex:single-item}
Consider a fair division instance with $n$ agents and one item
(which is either a good to everyone or a chore to everyone).
Then every allocation is EFX, EF1, APS, MMS, PROPx, GAPS, GMMS, PAPS, PMMS,
EEFX, EEF1, MXS, M1S, PROPm, and PROP1,
but not EF or PROP or GPROP or PPROP or EEF or MEFS.
\end{example}

\begin{lemma}
\label{cex:share-vs-envy-goods}
Consider a fair division instance $([n], [m], (v_i)_{i=1}^n, \eqEnt)$
with $n \ge 3$, $m = 2n-1$, and identical additive valuations,
where each item has value 1 to each agent.
Let $A$ be an allocation where agent $n$ gets $n$ goods,
and all other agents get 1 good each.
Then this allocation is APS+MMS+EEFX+PROPx, but not EF1.
\end{lemma}
\begin{proof}
By \cref{thm:impl:aps-to-pess,thm:impl:prop-to-aps},
$1 = \MMS_i \le \APS_i \le v([m])/n = 2 - 1/n$.
Hence, $APS_i = 1$, since the APS is the value of some bundle.
Hence, $A$ is APS+MMS+PROPx.
Agent $n$'s EEFX-certificate for $A$ is $A$ itself.
For $i \neq n$, agent $n$'s EEFX-certificate is $B$,
where $|B_i| = 1$ and $|B_j| = 2$ for $j \neq i$.
%
$A$ is not EF1 because the first $n-1$ agents EF1-envy agent $n$.
\end{proof}

\begin{lemma}
\label{cex:share-vs-envy-chores}
Consider a fair division instance $([n], [m], (v_i)_{i=1}^n, w)$
with $n \ge 3$, equal entitlements, identical additive disutilities,
and $m = n+1$ chores, each of disutility 1.
Let $A$ be an allocation where agents 1 and 2 get 2 chores each,
agent $n$ gets 0 chores, and the remaining agents get 1 chore each.
Then this allocation is APS+MMS+EEFX+PROPx, but not EF1.
\end{lemma}
\begin{proof}
$-2 = \MMS_i \le \APS_i \le v([m])/n = - 1 - 1/n$.
Hence, $APS_i = -2$, since the APS is the value of some bundle.
Hence, $A$ is APS+MMS+PROPx.

Agents $[n] \setminus [2]$ do not EFX-envy anyone in $A$.
For $i \in \{1, 2\}$, agent $i$'s epistemic-EFX-certificate can be obtained by
transferring a chore from agent $3-i$ to agent $n$.
Hence, $A$ is epistemic EFX.

$A$ is not EF1 because agents 1 and 2 EF1-envy agent $n$.
\end{proof}

\subsection{From EEF, MEFS, PROP}
\label{sec:cex-extra:from-eef-mefs-prop}

\begin{lemma}[EEF $\nfimplies$ EF1]
\label{cex:eef-not-ef1}
Let $0 \le 2a < b$. Let $f_1, f_2, f_3: 2^{[12]} \to \mathbb{R}_{>0}$ be additive sets functions:

\begin{tabular}{c|cccc|cccc|cccc}
& 1 & 2 & 3 & 4 & 5 & 6 & 7 & 8 & 9 & 10 & 11 & 12
\\ \hline $f_1$ & $a$ & $a$ & $b$ & $b$ & $b$ & $b$ & $b$ & $b$ & $a$ & $a$ & $a$ & $a$
\\ $f_2$ & $a$ & $a$ & $a$ & $a$ & $a$ & $a$ & $b$ & $b$ & $b$ & $b$ & $b$ & $b$
\\ $f_3$ & $b$ & $b$ & $b$ & $b$ & $a$ & $a$ & $a$ & $a$ & $a$ & $a$ & $b$ & $b$
\end{tabular}

Let $t \in \{-1, 1\}$ and let $\Ical \defeq ([3], [12], (v_i)_{i=1}^3, \eqEnt)$
be a fair division instance where $v_i \defeq tf_i$ for all $i \in [3]$.
Then allocation
$A \defeq ([4], [8] \setminus [4], [12] \setminus [8])$ is EEF+PROP but not EF1.
\end{lemma}
\begin{proof}
For $t = 1$, agent 1 EF1-envies agent 2 in $A$,
and for $t = -1$, agent 1 EF1-envies agent 3 in $A$.
$B = ([4], \{5, 6, 9, 10\}, \{7, 8, 11, 12\})$ is agent 1's EEF-certificate.
A similar argument holds for agents 2 and 3 too.
\end{proof}

\begin{example}[PROP $\nfimplies$ MEFS]
\label{cex:prop-not-mefs-goods}
Consider a fair division instance with 3 equally-entitled agents
having additive valuations over 3 goods:

\begin{tabular}{c|ccc}
& 1 & 2 & 3
\\ \hline $v_1$ & 10 & 20 & 30
\\ $v_2$ & 20 & 10 & 30
\\ $v_3$ & 10 & 20 & 30
\end{tabular}

Then the allocation $(\{2\}, \{1\}, \{3\})$ is PROP, but no allocation is MEFS
(every agent's minimum EF share is 30).
\end{example}

\begin{example}[PROP $\nfimplies$ MEFS]
\label{cex:prop-not-mefs-chores}
Consider a fair division instance with 3 equally-entitled agents
having additive disutilities over 3 chores:

\begin{tabular}{c|ccc}
& 1 & 2 & 3
\\ \hline $-v_1$ & 30 & 20 & 10
\\ $-v_2$ & 20 & 30 & 10
\\ $-v_3$ & 30 & 20 & 10
\end{tabular}

Then the allocation $(\{2\}, \{1\}, \{3\})$ is PROP, but no allocation is MEFS
(every agent's minimum EF share is $-10$).
\end{example}

\begin{lemma}[MEFS $\nfimplies$ EEF]
\label{cex:mefs-not-eef-goods}
Consider a fair division instance with 3 equally-entitled agents
having additive valuations over 6 goods:

\begin{tabular}{c|cccccc}
& 1 & 2 & 3 & 4 & 5 & 6
\\ \hline $v_1$ & 20 & 20 & 20 & 10 & 10 & 10
\\ $v_2$, $v_3$ & 20 & 10 & 10 &  1 &  1 &  1
\end{tabular}

Then the allocation $A \defeq (\{4, 5, 6\}, \{1\}, \{2, 3\})$ is MEFS, but no allocation is epistemic EF.
\end{lemma}
\begin{proof}
Agents 2 and 3 are envy-free in $A$.
Agent 1 has $B \defeq (\{1, 4\}, \{2, 5\}, \{3, 6\})$ as her MEFS-certificate for $A$.
Hence, $A$ is MEFS.

Suppose an epistemic EF allocation $X$ exists.
Let $Y^{(i)}$ be each agent $i$'s epistemic-EF-certificate.
For agent 2 to be envy-free in $Y^{(2)}$,
we require $Y^{(2)}_2 \supseteq \{1\}$ or $Y^{(2)}_2 \supseteq \{2, 3\}$.
Similarly, $Y^{(3)}_3 \supseteq \{1\}$ or $Y^{(3)}_3 \supseteq \{2, 3\}$.
Since $Y^{(i)}_i = X_i$ for all $i$, we get $X_2 \cup X_3 \supseteq \{1, 2, 3\}$.
Hence, $X_1 \subseteq \{4, 5, 6\}$.
But then no epistemic-EF-certificate exists for agent 1 for $X$,
contradicting our assumption that $X$ is epistemic EF.
Hence, no epistemic EF allocation exists.
\end{proof}

\begin{example}[MEFS $\nfimplies$ EEF]
\label{cex:mefs-not-eef-chores}
Consider a fair division instance with 3 equally-entitled agents
having additive disutilities over 6 chores:

\begin{tabular}{c|cccccc}
& 1 & 2 & 3 & 4 & 5 & 6
\\ \hline $-v_1$ & 20 & 20 & 20 & 10 & 10 & 10
\\ $-v_2$, $-v_3$ & 20 & 10 & 10 & 10 & 10 & 10
\end{tabular}

Then the allocation $A \defeq (\{4, 5, 6\}, \{1\}, \{2, 3\})$ is MEFS
(agents 2 and 3 are EF, agent 1's MEFS-certificate is $(\{1, 4\}, \{2, 5\}, \{3, 6\})$).
Agent 1 is not epistemic-EF-satisfied by $A$.
\end{example}

\begin{lemma}[MEFS $\nfimplies$ EEF1]
\label{cex:mefs-not-eef1-chores}
Consider a fair division instance with 3 equally-entitled agents
having additive disutilities over 12 chores.
$v_1(1) = v(2) = v(3) = 70$ and $v(c) = 10$ for all $c \in [12] \setminus [3]$.
Agents 2 and 3 have disutility 10 for each chore.
Then $A \defeq ([12] \setminus [3], [2], \{3\})$ is a MEFS+PROP allocation
where agent 1 is not EEF1-satisfied.
\end{lemma}
\begin{proof}
$\PROP_1 = -100$ and $\PROP_2 = \PROP_3 = -40$.
$\MEFS_1 \le -100$ because of the allocation $(\{1, 4, 5, 6\}, \{2, 7, 8, 9\}, \{3, 10, 11, 12\})$.
$\MEFS_i \le -40$ for $i \in \{2, 3\}$ because of the allocation
$([4], [8] \setminus [4], [12] \setminus [8])$.
Agent 1 has disutility $90$ in $A$, so $A$ is MEFS-fair and PROP-fair to agent 1.
Agents 2 and 3 have disutility at most $20$ in $A$, so $A$ is MEFS-fair and PROP-fair to them.

Agent 1 is not EEF1-satisfied by $A$, since in any EEF1-certificate $B$,
some agent $j \in \{2, 3\}$ receives at most one chore of value $70$,
and agent 1 would EF1-envy $j$.
\end{proof}

\subsection{Two Equally-Entitled Agents}
\label{sec:cex-extra:2-eqEnt}

\begin{example}[EFX $\nfimplies$ MMS]
\label{cex:efx-not-mms}
Let $t \in \{-1, 1\}$.
Consider a fair division instance with 2 equally-entitled agents having
an identical additive valuation function $v$ over 5 items.
$v(1) = v(2) = 3t$ and $v(3) = v(4) = v(5) = 2t$.
Then allocation $A \defeq (\{1, 3\}, \{2, 4, 5\})$ is EFX.
The MMS is $6t$, since $P = (\{3t, 3t\}, \{2t, 2t, 2t\})$ is an MMS partition.
But in $A$, some agent doesn't get her MMS.
\end{example}

\begin{example}[EF1 $\nfimplies$ PROPX or MXS]
\label{cex:ef1-not-propx-mxs}
Let $t \in \{-1, 1\}$.
Consider a fair division instance with 2 equally-entitled agents
having an identical additive valuation function $v$ over 5 items,
where $v(1) = v(2) = 4t$ and $v(3) = v(4) = v(5) = t$.
Then allocation $A \defeq (\{1\}, [5] \setminus \{1\})$ is EF1 but not PROPx and not MXS.
\end{example}

\begin{lemma}[PROPx $\nfimplies$ M1S]
\label{cex:propx-not-m1s}
Let $t \in \{-1, 1\}$ and $0 < \eps < 1/2$.
Consider a fair division instance with 2 equally-entitled agents having
an identical additive valuation function $v$ over 4 items.
Let $v(4) = (1+2\eps)t$ and $v(j) = t$ for $j \in [3]$.
Then allocation $A \defeq (\{4\}, [3])$ is PROPx but not M1S.
\end{lemma}
\begin{proof}
$v([m])/2 = (2+\eps)t$, so $A$ is PROPx.
%
For $t = 1$, in any allocation $B$ where agent 1 doesn't EF1-envy agent 2, she must have at least 2 goods.
But $v(A_1) = 1+2\eps$, so agent 1 doesn't have an M1S-certificate for $A$. Hence, $A$ is not M1S.
%
For $t = -1$, in any allocation $B$ where agent 2 doesn't EF1-envy agent 1, she must have at most 2 chores.
But $v(A_2) = -3$, so agent 1 doesn't have an M1S-certificate for $A$. Hence, $A$ is not M1S.
\end{proof}

\begin{example}[MXS $\nfimplies$ PROPx for $n=2$]
\label{cex:mxs-not-propx-n2}
Let $t \in \{-1, 1\}$.
Consider a fair division instance with 2 equally-entitled agents
having identical additive valuations over 7 items:
the first 2 items of value $4t$ and the last 5 items of value $t$.
Then the allocation $A = (\{1, 3\}, \{2, 4, 5, 6, 7\})$ is not PROPx or EFX,
but it is MXS because the agents have $([7] \setminus [2], [2])$
and $([2], [7] \setminus [2])$ as their MXS-certificates for $A$.
\end{example}

\begin{lemma}[M1S $\nfimplies$ PROP1]
\label{cex:m1s-not-prop1}
Consider a fair division instance with 2 equally-entitled agents
having an identical additive valuation function $v$ over 9 items.
Let $t \in \{-1, 1\}$ and $v(9) = 4t$ and $v(j) = t$ for $j \in [8]$.
Then allocation $A \defeq (\{9\}, [8])$ is M1S but not PROP1.
\end{lemma}
\begin{proof}
$v([9])/2 = 6t$. Let $B \defeq ([4], [9] \setminus [4])$.
%
For $t = 1$ (goods), $B$ is agent 1's M1S-certificate for $A$,
but agent 1 is not PROP1-satisfied by $A$.
%
For $t = -1$ (chores), $B$ is agent 2's M1S-certificate for $A$,
but agent 2 is not PROP1-satisfied by $A$.
\end{proof}

\subsection{Three Equally-Entitled Agents}
\label{sec:cex-extra:3-eqEnt}

\begin{example}[GAPS $\nfimplies$ PROPx]
\label{cex:gaps-not-propx}
Consider a fair division instance with 3 equally-entitled agents
having identical additive valuations. There are 2 goods of values 50 and 10.
In every allocation, some agent doesn't get any good, and that agent is not PROPx-satisfied.
The allocation where the first agent gets the good of value 5
and the second agent gets the good of value 1 is a groupwise APS allocation
(set the price of the goods to $1.1$ and $0.9$).
\end{example}

\begin{lemma}[APS $>$ MMS]
\label{thm:aps-gt-mms}
Let $t \in \{-1, 1\}$.
Consider a fair division instance with 3 equally-entitled agents
having identical additive valuations over 15 items. The items' values are
$65t$, $31t$, $31t$, $31t$, $23t$, $23t$, $23t$, $17t$, $11t$, $7t$, $7t$, $7t$, $5t$, $5t$, $5t$.
Then the AnyPrice share is at least $97t$, the proportional share is $97t$,
and the maximin share is less than $97t$.
\end{lemma}
\begin{proof}
For $t = 1$, this follows from Lemma C.1 of \cite{babaioff2023fair}.
For $t = -1$, a similar argument tells us that the AnyPrice share is at least $-97$.
If the maximin share is at least $-97$, then there must exist a partition $P$ of the chores
where each bundle has disutility 97. But then $P$ would prove that the maximin share
in the corresponding goods instance is at least 97, which is a contradiction.
Hence, for $t = -1$, the maximin share is less than $-97$.
\end{proof}

\begin{example}[GMMS $\nfimplies$ APS]
\label{cex:gmms-not-aps}
For the fair division instance in \cref{thm:aps-gt-mms},
the leximin allocation is GMMS (since on restricting to any subset of agents,
the resulting allocation is still leximin, and is therefore MMS).
However, no APS allocation exists, because APS $>$ MMS,
and the minimum value across all bundles is at most the MMS.
\end{example}

\begin{example}[PMMS $\nfimplies$ MMS, Example 4.4 of \cite{caragiannis2019unreasonable}]
\label{cex:pmms-not-mms}
Let $t \in \{-1, 1\}$.
Consider a fair division instance with 3 equally-entitled agents
having an identical valuation function $v$ over 7 items where
$v(1) = 6t$, $v(2) = 4t$, $v(3) = v(4) = 3t$, $v(5) = v(6) = 2t$, $v(7) = t$.
Each agent's maximin share is $7t$ ($(\{1, 7\}, \{2, 3\}, \{4, 5, 6\})$ is a maximin partition).
Allocation $(\{1\}, \{3, 4, 5\}, \{2, 6, 7\})$ is PMMS but not MMS.
\end{example}

\begin{example}[APS $\nfimplies$ PROPm]
\label{cex:aps-not-propm}
Consider a fair division instance with 3 equally-entitled agents
having an identical additive valuation function $v$ over 6 goods:
$v(1) = 60$, $v(2) = 30$, and $v(3) = v(4) = v(5) = v(6) = 10$.
The allocation $A \defeq (\{2\}, \{3, 4, 5\}, \{1, 6\})$ is APS+MMS, since the MMS is 30,
and the APS is at most 30 because of the price vector $(4, 3, 1, 1, 1, 1)$.
However, $A$ is not PROPm-fair to agent 1, because the proportional share is $130/4 > 40$.
\end{example}

\begin{example}[APS $\nfimplies$ PROP1]
\label{cex:aps-not-prop1-chores}
Consider a fair division instance with 3 equally-entitled agents
having an identical additive valuation function $v$ over 6 chores:
the first chore has disutility 18 (large chore) and the remaining chores have disutility 3 each (small chores).
Then $X \defeq ([6] \setminus \{1\}, \{1\}, \emptyset)$ is MMS+APS,
since the MMS is $-18$, and the APS is at most $-18$ due to the price vector $(1, 0, 0, 0, 0, 0)$.
$X$ is not PROP1-fair to agent 1, since the proportional share is $-11$,
and agent 1's disutility in $X$ after removing any chore is $12$.
$X$ is not EEF1-fair to agent 1 because even after redistributing chores among the remaining agents,
someone will always have no chores.
\end{example}

\begin{lemma}[PROPm doesn't exist for mixed manna]
\label{cex:propm-mixed-manna}
Consider a fair division instance $([3], [6], (v_i)_{i=1}^3, \eqEnt)$
where agents have identical additive valuations, and the items have values
$(-3, -3, -3, -3, -3, 3\eps)$, where $0 < \eps < 1/2$.
Then there exists an EFX+GMMS+GAPS allocation but no PROPm allocation.
\end{lemma}
\begin{proof}
The proportional share is $v([m])/3 = -5 + \eps$.
\WLoG, assume agent 1 receives the most number of chores,
and agent 3 receives the least number of chores.
Then agent 1 has at least 2 chores, and agent 3 has at most 1 chore.

\textbf{Case 1}: agent 1 receives at least 3 chores.
\\ Then even after removing one of her chores, and even if she receives the good,
her value for her bundle is at most $-6 + 3\eps < -5 + \eps$.
Hence, she is not PROPm-satisfied.

\textbf{Case 2}: agent 1 receives 2 chores.
\\ Then agent 2 also receives 2 chores, and agent 3 receives 1 chore.
Assume without loss of generality that agent 2 receives the good.
Then this allocation is EFX and groupwise MMS.
On setting the price of each chore to $-3$ and the price of the good to 3,
we get that the allocation is groupwise APS.
However, if agent 1 adds the good to her bundle, her value becomes $-6 + 3\eps < -5 + \eps$.
Hence, she is not PROPm-satisfied.
\end{proof}

\subsection{Unequal Entitlements}
\label{sec:cex-extra:uneqEnt}

\begin{lemma}[PROP1+M1S is infeasible]
\label{cex:prop1-plus-m1s-ue}
Consider a fair division instance $\Ical \defeq ([3], [7], (v_i)_{i=1}^3, w)$,
where the entitlement vector is $w \defeq (7/12, 5/24, 5/24)$,
the agents have identical additive valuations, and each good has value 1. Then
\begin{tightenum}
\item $\APS_1 = 4$ and $\APS_2 = \APS_3 = 1$.
\item $X$ is APS $\iff$ $X$ is groupwise-APS (GAPS) $\iff$ $X$ is PROP1.
\item $\WMMS_1 = 3$ and $\WMMS_2 = \WMMS_3 = 15/14$.
\item $X$ is WMMS $\iff$ $X$ is groupwise-WMMS (GWMMS) $\iff$ $X$ is EFX $\iff$ $X$ is M1S.
\end{tightenum}
Therefore,
\begin{tightenum}
\item M1S+PROP1 is infeasible for this instance.
\item GWMMS+EFX doesn't imply PROP1.
\item GAPS doesn't imply M1S.
\end{tightenum}
\end{lemma}
\begin{proof}
By \cref{thm:impl:tribool:aps}, $\APS_1 = \floor{\frac{7 \times 7}{12}} = 4$
and $\APS_2 = \APS_3 = \floor{\frac{5 \times 7}{24}} = 1$.
By \cref{thm:impl:tribool:aps,thm:impl:tribool:prop1},
an allocation is APS iff it is PROP1.

Any GAPS allocation is also APS by definition.
We will now show that any APS allocation is also GAPS.
Formally, let $A$ be an APS allocation for $\Ical$.
The cardinality vector of $A$, i.e., $c \defeq (|A_1|, |A_2|, |A_3|)$,
can have three possible values: $(5, 1, 1)$, $(4, 2, 1)$, $(4, 1, 2)$.
For every possible value of $c$ and $S \subseteq [3]$,
we show that $(\Icalhat, \Ahat) \defeq \restrict(\Ical, A, S)$ is APS
(c.f.~\cref{defn:restricting}).
\begin{tightenum}
\item $c = (5, 1, 1)$ and $S = \{1, 2\}$:
    $\Icalhat$ has 6 goods and entitlement vector $(14/19, 5/19)$.
    $\APS_1 = \floor{\frac{14 \times 6}{19}} = 4$ and $\APS_2 = \floor{\frac{5 \times 6}{19}} = 1$.
    Hence, $\Ahat$ is APS for $\Icalhat$.
\item $c = (5, 1, 1)$ and $S = \{1, 3\}$:
    Similar to the $S = \{1, 2\}$ case.
\item $c = (5, 1, 1)$ and $S = \{2, 3\}$:
    $\Icalhat$ has 2 goods and entitlement vector $(1/2, 1/2)$.
    $\APS_3 = \APS_3 = 1$, so $\Ahat$ is APS for $\Icalhat$.
\item $c = (4, 2, 1)$ and $S = \{1, 2\}$:
    $\Icalhat$ has 6 goods and entitlement vector $(14/19, 5/19)$.
    $\APS_1 = \floor{\frac{14 \times 6}{19}} = 4$ and $\APS_2 = \floor{\frac{5 \times 6}{19}} = 1$.
    Hence, $\Ahat$ is APS for $\Icalhat$.
\item $c = (4, 2, 1)$ and $S = \{1, 3\}$:
    $\Icalhat$ has 5 goods and entitlement vector $(14/19, 5/19)$.
    $\APS_1 = \floor{\frac{14 \times 5}{19}} = 3$ and $\APS_2 = \floor{\frac{5 \times 5}{19}} = 1$.
    Hence, $\Ahat$ is APS for $\Icalhat$.
\item $c = (4, 2, 1)$ and $S = \{2, 3\}$:
    $\Icalhat$ has 3 goods and entitlement vector $(1/2, 1/2)$.
    $\APS_2 = \APS_3 = \floor{\frac{1 \times 3}{2}} = 1$.
    Hence, $\Ahat$ is APS for $\Icalhat$.
\item $c = (4, 1, 2)$:
    Similar to the $c = (4, 2, 1)$ case.
\end{tightenum}
Hence, any APS allocation for $\Ical$ is also GAPS.

For any allocation $X$, define
\[ f(X) \defeq \min_{j=1}^3 \frac{|X_j|}{w_j}. \]
Then $\WMMS_i = w_i\max_X f(X)$ for all $i \in [3]$.
If $|X_1| \le 2$, then $f(X) \le |X_1|/w_1 \le 24/7 = 3 + 3/7$.
If $|X_2| \le 1$ or $|X_3| \le 1$, $f(X) \le 24/5 = 4 + 4/5$.
Otherwise, $|X_1| = 3$ and $|X_2| = |X_3| = 2$,
so $f(X) = \min(3 \times 12/7, 2 \times 24/5) = 36/7 = 5 + 1/6$.
Hence, $\max_X f(X) = 36/7$, so $\WMMS_1 = 3$ and $\WMMS_2 = \WMMS_3 = 15/14$.
So, an allocation is WMMS iff it has cardinality vector $(3, 2, 2)$.

Any GWMMS allocation is also WMMS by definition. We now prove the converse.
For every possible value of $S \subseteq [3]$,
we show that $(\Icalhat, \Ahat) \defeq \restrict(\Ical, A, S)$ is WMMS
(c.f.~\cref{defn:restricting}).
\begin{tightenum}
\item $S = \{1, 2\}$:
    $\Icalhat$ has 5 goods and entitlement vector $(14/19, 5/19)$.
    If $|X_1| \le 2$, then $f(X) \le |X_1|/w_1 \le 38/14 = 2 + 10/14$.
    If $|X_2| \le 1$, then $f(X) \le |X_2|/w_2 \le 19/5 = 3 + 10/14$.
    Otherwise, $|X_1| = 3$ and $|X_2| = 2$, so $f(X) = \min(3/w_1, 2/w_2) = \min(57/14, 38/5) = 57/14 = 4 + 1/14$.
    Hence, $\WMMS_1 = 3$ and $\WMMS_2 = \WMMS_3 = 57/14 \times 5/19 = 1 + 19/266$.
    Hence, $\Ahat$ is WMMS for $\Icalhat$.
\item $S = \{1, 3\}$:
    Similar to the $S = \{1, 2\}$ case.
\item $S = \{2, 3\}$:
    $\Icalhat$ has 4 goods and entitlement vector $(1/2, 1/2)$.
    Then $\WMMS_1 = \WMMS_2 = 2$.
    Hence, $\Ahat$ is WMMS for $\Icalhat$.
\end{tightenum}
Hence, any WMMS allocation for $\Ical$ is also GWMMS.

Any GWMMS allocation is EFX by \cref{thm:impl:mms-to-efx-n2},
and any EFX allocation is M1S by \cref{thm:impl:efx-to-ef1}.
We will now show that any M1S allocation is WMMS.

Let $X$ be an M1S allocation.
Let $A$ be agent 1's M1S certificate.
If $|A_1| \le 2$, then $|A_j| \ge 3$ for some $j \in \{2, 3\}$.
Since agent 1 has higher entitlement, she would EF1-envy agent $j$,
which is a contradiction. Hence, $|X_1| \ge |A_1| \ge 3$.

Let $B$ be agent 2's M1S certificate.
Suppose $|B_2| \le 1$. Since $2$ doesn't EF1-envy $3$, we get $|B_3| \le 2$.
Since $2$ doesn't EF1-envy $1$, we get
\[ \frac{|B_1|-1}{w_1} \le \frac{|B_2|}{w_2} \iff |B_1| \le 1 + \frac{w_1}{w_2} = 3 + \frac{4}{5}. \]
Hence, $|B_1| + |B_2| + |B_3| \le 3 + 1 + 2 = 6$, which is a contradiction.
Hence, $|X_2| \ge |B_2| \ge 2$.
Similarly, we can prove that $|X_3| \ge 2$.

Hence, $|X_i| \ge \WMMS_i$ for all $i$, so $X$ is WMMS.
This proves that any M1S allocation is WMMS.
\end{proof}

\cite{chakraborty2021weighted} also proves that EF1+PROP1 allocations may not exist for unequal entitlements,
We use a different counterexample in \cref{cex:prop1-plus-m1s-ue},
which allows us to also prove other non-implications.

\begin{example}[PROP1 $\nfimplies$ M1S]
\label{cex:prop1-not-m1s-n2}
Consider a fair division instance with 2 agents having identical additive valuations.
Let $t \in \{-1, 1\}$. Let there be 2 items, each of value $t$.
Let the entitlement vector be $(2/3, 1/3)$.
Let $A$ be an allocation where the first agent gets both items.
Then $A$ is PROP1 but not M1S.
\end{example}

\subsection{Non-Additive Valuations}
\label{sec:cex-extra:non-add}

\begin{lemma}[EF $\nfimplies$ PROP]
\label{cex:ef-not-prop-supmod}
Consider a fair division instance with 2 agents and 4 goods.
The agents have identical valuations and equal entitlements.
Let $a, b \in \mathbb{R}_{\ge 0}$ such that $3a < b$.
The valuation function $v$ is given by
\[ v(S) \defeq \begin{cases}
|S|a & \textrm{ if } |S| \le 3
\\ 3a + b & \textrm{ if } |S| = 4
\end{cases}. \]
Then $v$ is supermodular, no PROP1 allocation exists,
and if $a > 0$, then no PROPm allocation exists.
However, an allocation where each agent gets 2 goods is EF+APS+MMS.
\end{lemma}
\begin{proof}
For any $g \in [4]$ and $S \subseteq [4] \setminus \{g\}$, we have
\[ v(g \mid S) = \begin{cases}
a & \textrm{ if } |S| \le 2
\\ b & \textrm{ if } |S| = 3
\end{cases}. \]
Hence, $v$ is supermodular.

The proportional share is $(3a + b)/2$.
In any allocation, some agent gets at most 2 goods,
and even if she is given an additional good, her valuation is $3a < (3a+b)/2$.
Hence, no allocation is PROP1, and no allocation is PROPm if $a > 0$.
(When $a = 0$, every allocation is PROPm.)

$v(A_1) = v(A_2) = 2a$, so $A$ is EF.
It is easy to check that the MMS is $2a$.

If we set the price of each good to 1, then at most 2 goods are affordable.
Hence, APS is at most $2a$.
Moreover, for any price vector, the cheapest 2 goods are affordable, and their total valuation is $2a$.
Hence, APS is at least $2a$.
\end{proof}

We first define a function in \cref{cex:ud-submod-canc} by perturbing a unit-demand function.
We then use it to show that a PROP allocation may not be M1S.

\begin{lemma}
\label{cex:ud-submod-canc}
Let $0 \le \eps < 1/6$ and $v: 2^{[3]} \to \mathbb{R}_{\ge 0}$, where
\[ v(S) := \begin{cases}
4 + 2\eps|S| & \textrm{ if } 1 \in S \textrm{ or } 2 \in S
\\ 3 + 2\eps|S| & \textrm{ if } S = \{3\}
\\ 0 & \textrm{ otherwise}
\end{cases}. \]
Then $v$ is submodular and cancelable.
\end{lemma}
\begin{proof}
\[ v(g \mid S) = 2\eps + \begin{cases}
0 & \textrm{ if } 1 \in S \textrm{ or } 2 \in S
\\ 1 & \textrm{ if } S = \{3\}
\\ v(\{g\}) & \textrm{ if } S = \emptyset
\end{cases}. \]
We can see that adding elements to $S$ never increases $v(g \mid S)$.
Hence, $v$ is submodular.

One can confirm that $v$ is cancelable by painstakingly applying the definition of cancelable valuations
(c.f.~\cref{sec:settings-extra}).
\end{proof}

\begin{lemma}[PROP $\nfimplies$ M1S]
\label{cex:prop-not-m1s-submod}
Consider a fair division instance with 2 equally-entitled agents
having an identical valuation function $v$ over 3 goods as defined in \cref{cex:ud-submod-canc}.
Then the allocation $A = (\{1, 2\}, \{3\})$ is PROP but not M1S.
\end{lemma}
\begin{proof}
$A$ is PROP since $v(A_1) = 4+6\eps$, $v(A_2) = 3+2\eps$, and the PROP share is $2+3\eps$.
Suppose $A$ is M1S and agent 2's M1S certificate for $A$ is $B$.
Then $v(B_2) \le v(A_2) = 3 + 2\eps$, so $B_2 = \{3\}$. Hence, $B_1 = \{1, 2\}$.
However, $\min_{g \in B_1} v(B_1 \setminus \{g\}) = 4 + 6\eps > 3 + 2\eps = v(B_2)$.
Hence, agent 2 is not EF1-satisfied by $B$, which contradicts the fact that $B$
is agent 2's M1S certificate for $A$. Hence, $A$ is not M1S.
\end{proof}
